% --------------------------------------------------------------------------
% Template for ICAD-2023 Sonic Tilt paper; to be used with:
%          icad2023.sty  - ICAD 2023 LaTeX style file, and 2023
%          IEEEbtran.bst - IEEE bibliography style file.
%
% --------------------------------------------------------------------------
\UseRawInputEncoding
\documentclass[a4paper,10pt,oneside]{article}
\usepackage[utf8]{inputenc}
\usepackage{icad2023,amsmath,epsfig,times,url,hyperref}
\usepackage[T1]{fontenc}

% Example definitions.
% --------------------
\def\defeqn{\stackrel{\triangle}{=}}
\newcommand{\symvec}[1]{{\mbox{\boldmath $#1$}}}
\newcommand{\symmat}[1]{{\mbox{\boldmath $#1$}}}

% Title.
% --------------------
\title{AUTHOR GUIDELINES FOR ICAD 2023 SONIC TILT COMPETITION SUBMISSIONS}
% *** IMPORTANT ***
% *** PLEASE LEAVE AUTHOR INFORMATION BLANK UNTIL FINAL CAMERA-READY SUBMISSION *** 

% IF ONE AUTHOR
%\name{Jyri Huopaniemi} 
%\address{Nokia Research Center \\ 
%Speech and Audio Systems Laboratory \\ 
%P.O.Box 407, FIN-00045 Nokia Group, Finland \\ 
%{\tt jyri.huopaniemi@nokia.com}} 
%

% IF TWO AUTHORS
\twoauthors{Author 1} {Organization \\ Address Line 2 \\ City, Country  \\ {\tt author1@example.com}}
{Author 2} {Organization \\ Address Line 2 \\ City, Country  \\ {\tt author2@example.com}}

\begin{document}
\ninept
\maketitle
\begin{sloppy}
\begin{abstract}
This is the template file for manuscript submissions to the Sonic Tilt Competition at the 2023 International Conference on Auditory Display, which will be held June 26 to July 1, 2023 at the Linköping University, Norrköping, Sweden. This template has been generated from the ICAD2023 template and aims at producing a Sonic Tilt Competition paper collection that is not a part of the official ICAD2023 proceedings. Please use \LaTeX\ when preparing your submission. All questions concerning ICAD2023-Sonic Tilt Competition submissions should be addressed to the Sonic Tilt chair (\href{mailto:sonictilt@icad2023.icad.org}{sonictilt@icad2023.icad.org}). The template is available in electronic form at the website: \href{https://icad2023.icad.org/}{https://icad2023.icad.org/} 
Further information on the Sonic Tilt Competition can be found on the \href{https://icad2023.icad.org/sonic-tilt/}{ICAD2023 website} and on the  \href{https://tiltification.uni-bremen.de/competition.html}{Tiltification website}. Please give your app a catchy name (ideally not ``Sonic Tilt'' or ``Tiltification'') and don't forget to include explanatory graphics and provide links to helpful videos, websites, and, most importantly, the APK file that you compiled.
\end{abstract}
\section{Link to APK File}
\label{sec:link}
Don't forget to provide us with the built app from your flutter folder. The APK file of the original Tiltification app can be found under \url{https://tiltification.uni-bremen.de/assets/download/tiltification-app.apk}.
\section{Introduction}
\label{sec:intro}

The introduction should explain the motivation, concept, or reasoning of your sonification approach, and/or your progress of mind-storming or source of inspiration. It can be based on existing methods, your personal or professional experience. The question is: How shall your sonification convince the masses to use sonification? This is the main criterion for the app. Have you found the ideal compromise between precision and pleasantness of the sound? Is it fun to learn to interpret the sonification? Is it a musical experience, or a thrilling acoustic expedition? Are people already familiar with the sound, or is it new and exciting? Your app name and icon may reflect your sonification design idea. Your submission should have a length between $2$ and $4$ pages, depending on the number and size of graphics.

\section{Your Sonification}
\label{sec:format}
Please explain how your sound will enable app users to level their phones, guided by your sonification. You should explain the sound in general as well as how the sound changes, depending on the tilt angles. How do users know whether they need to lift the left side of the phone, or the right side? And how do they know by how far they need to lift it? How do they know whether they have to tilt it towards them, or away from them? And how far? And how do combinations sound? Does your sound inform them whether they are on track or not? How do they know when their phone is leveled? Ideally, non-experts should be able to understand and use your sonification after your explanation. Explaining the mapping principle in colloquial terms is certainly a way to explain your sonification. Illustrations and links to videos may help a lot. For example, you can make a screenshot of the app and add some explanatory text, icons, plots or graphics, as illustrated in Fig. \ref{pic:screenshot}. You can find a screenshot of the app, called \verb!blank.png!,  in the \verb!pics! folder of the \LaTeX{} template.

\begin{figure}[ht]
\centerline{\includegraphics[width=83mm]{icad2022_LATEXtemplate/pics/screenshot}}
	\caption{{\it Blending colored graphics, text and symbols with a screenshot of your app makes it easy to understand, especially in combination with your written explanation.}}
	\label{pic:screenshot}
\end{figure}

I used the \href{https://play.google.com/store/apps/details?id=com.blogspot.byterevapps.lollipopscreenrecorder&hl=de&gl=US}{ADV screen recorder} to record video and sound of my smartphone while using the app, as you can see on  \url{https://youtu.be/Id_BkvpB07c}. We also shot some short clips of people using our app, as in \url{https://youtu.be/STy6VRn0s3o}.

Don't forget to upload your app to a server and provide a permanent download link to the APK file. By participating in the competition, you agree that ICAD may archive your app and your submitted description. Naturally, you will be credited for your app.

\section{Implementation}
\label{sec:pagestyle}
Here, you can describe the technical implementation of your sonification. Do you use FM-synthesis \cite{chown}, the Karplus-Strong \cite{karplus} algorithm, or the Impulse Pattern Formulation \cite{ipfchaos}? Have you implemented a sonification that has already been published? Then you should cite the respective source. You can also link to your PD file or show a screenshot. This section is for geeks and freaks who may want to understand technical details or even use your sonification themselves. Of course, your implementation may also stay a company secret if you like. 

\section{Evaluation}
\label{sec:typestyle}

A proper evaluation of your sonification is not necessary to participate in the Sonic Tilt competition. A jury will test your sonification and read your app description to select the competition winner. Of course, you can evaluate your approach using PAMPAS \cite{pampas}, BUZZ \cite{buzz}, NASA-TLX \cite{nasaoriginal} and other techniques. Personally, I carried out a short performance experiment with Tiltification \cite{isontable}, whose results may serve as a benchmark for your sonification evaluation.

If you carry out an experimental evaluation of your sonification, you should submit a dedicated scientific paper in addition to your submission for the Sonic Tilt Competition. Note, however, that the submission deadline for ICAD2023 is much earlier than the submission deadline for the Sonic Tilt Competition. 

\section{Submission}
It will be possible to submit your app description via easy chair, following the link that will be provided on the ICAD2023 website. Make sure to submit it as a contribution to the Sonic Tilt Competition.

The submission deadline is 12th May 2023.

During the ICAD2023, you will have the chance to give a 5-minute teaser/demonstration of your app. To participate in the competition only, it is not necessary to register for the ICAD2023 conference and pay the fees. You can attend the Sonic Tilt Competition session online and on-site. However, participants of the competition can only give their presentation and watch the presentation of their competitors. To attend the rest of the conference, you have to register and pay the registration fees.

\section{Acknowledgment}
\label{sec:ack}
Big thanks to Stephen Barrass who had the idea for the Sonic Tilt Competition and who already provided feedback on installing and building the app. I also thank Niklas Rönnberg who gave the Sonic Tilt Competition a spot in the ICAD2023 conference.

As you are using our open source project \href{https://github.com/Tiltification/sonic-tilt}{Sonic Tilt}, please cite our paper that introduces Tiltification, the original version of Sonic Tilt \cite{tilt}. In the spirit of the International Community for Auditory Display, I recommend all participants to use their app as a demonstrator for science communication and teaching, and as an advertisement for sonification. In \cite{recommendations} we give recommendations on how to develop, distribute and market sonification apps.
% -------------------------------------------------------------------------
% Either list references using the bibliography style file IEEEtran.bst
\bibliographystyle{IEEEtran}
\bibliography{refs2023}
%
% or list them by yourself
% \begin{thebibliography}{9}
% 
% \bibitem{icad2015web}
%   \url{http://www.icad.org}.
%
%\bibitem[1]{icad1} A.~Bee, C.D.~Player, and X.~Lastname, ``A correct citation,'' in {\it Proc. of the 1st Int. Conf. (IC)}, Helsinki, Finland, June 2001, pp. 1119-1134.  
%\bibitem[2]{icad2} E.~Zwicker and H.~Fastl, {\it Psychoacoustics: Facts and Models}, Springer-Verlag, Heidelberg, Germany, 1990.
%\bibitem[3]{icad3} M.R.~Smith, ``A good journal article,'' {\it J. Acoust. Soc. Am.}, vol. 110, no. 3, pp. 1598--1608, Mar. 2001.
% 
% \end{thebibliography}

\end{sloppy}
\end{document}
